% Part: first-order-logic
% Chapter: model-theory
% Section: nonstandard-arithmetic

\documentclass[../../../include/open-logic-section]{subfiles}

\begin{document}

\olfileid[id]{mod}{bas}{nsa}
\section{Model Aritmetika Tak Standar}

\begin{defn}
Misalkan $\Lang{L_N}$ adalah !!{language} aritmetika, yang terdiri atas
!!{constant} $\Obj{0}$, !!{predicate} dua-tempat $<$, !!{function}
satu-tempat $\prime$, serta !!{function}s dua-tempat $+$ dan $\times$.
\begin{enumerate}
\item \emph{Model standar} aritmetika adalah !!{structure} $\Struct{N}$
  untuk $\Lang{L_N}$ dengan $\Nat = \{0, 1, 2, \dots\}$, yang
  menafsirkan $\Obj{0}$ sebagai $0$, $<$ sebagai relasi kurang-dari
  pada~$\Nat$, serta $\prime$, $+$, dan $\times$ masing-masing sebagai
  operasi penerus, penjumlahan, dan perkalian pada~$\Nat$.
\item \emph{Aritmetika benar} adalah teori $\Theory{N}$.
\end{enumerate}
\end{defn}

Ketika bekerja dalam $\Lang{L_N}$, setiap suku berbentuk
$\Obj{0}^{\prime\dots\prime}$, dengan $n$ kali penerapan fungsi penerus
pada $\Obj{0}$, kita singkat sebagai~$\num{n}$.

\begin{defn}
Suatu !!{structure}~$\Struct{M}$ untuk $\Lang{L_N}$ bersifat
\emph{standar} jika dan hanya jika $\Struct{N} \iso \Struct{M}$.
\end{defn}

\begin{thm}\ollabel{thm:non-std}
Terdapat model-model !!{enumerable} tak standar dari aritmetika benar.
\end{thm}

\begin{proof}
Perluas $\Lang{L_N}$ dengan memperkenalkan !!{constant} baru~$c$, dan
perhatikan teori
\[
\Theory{N} \cup \Setabs{\num{n} < c}{n \in \Nat}.
\]
Teori tersebut dapat dipenuhi secara berhingga, sehingga menurut kekompakan teori
itu mempunyai suatu model $\Struct{M}$, yang menurut teorema
L\"owenheim--Skolem Menurun dapat dipilih !!{enumerable}. Jika
$\Domain{M}$ adalah ranah $\Struct{M}$, misalkan $\Struct{M}$ menafsirkan
konstanta-konstanta nonlogis dari $\Lang{L}$ sebagai $\mathbf{z} =
\Assign{\Obj{0}}{M} \in \Domain{M}$, ${\prec} = \Assign{<}{M} \subseteq
M^2$, $* = \Assign{\prime}{M} \colon \Domain{M} \to \Domain{M}$, serta
$\oplus = \Assign{+}{M}, \otimes = \Assign{\times}{M} \colon
\Domain{M}^2 \to \Domain{M}$. Untuk setiap $x \in \Domain{M}$, kita
menulis $x^*$ bagi unsur $\Domain{M}$ yang diperoleh dari $x$ dengan
menerapkan~$*$.

Sekarang, seandainya $h$ merupakan suatu isomorfisme antara $\Struct{N}$
dan $\Struct{M}$, akan ada $n \in \Nat$ sedemikian sehingga $h(n) =
\Assign{c}{M}$. Jadi, ambil sebarang penetapan $s$ dalam $\Struct{N}$
sedemikian sehingga $s(x) = n$. Maka $\Sat{N}{\eq[\num{n}][x]}[s]$;
menurut bukti \olref[iso]{thm:isom}, juga
$\Sat{M}{\eq[\num{n}][x]}[h \circ s]$, sehingga $\Assign{c}{M} =
\mathbf{z}^{*\cdots *}$ (dengan $*$ diiterasikan $n$ kali). Namun, hal ini
mustahil karena menurut asumsi $\Sat{M}{\num{n} < c}$ dan $\prec$
irefleksif. Jadi, $\Struct{M}$ tak standar.
\end{proof}

\begin{prob}
Suatu relasi $R$ pada himpunan $X$ \emph{berdasar-baik} jika dan hanya jika
tidak terdapat rantai menurun tak hingga dalam~$R$, yakni jika tidak ada
$x_0$, $x_1$, $x_2$,~\dots dalam $X$ sedemikian sehingga
$\dots x_2Rx_1Rx_0$. Dengan mengasumsikan teori himpunan
Zermelo--Fraenkel~$ZF$ konsisten, tunjukkan bahwa terdapat model-model $ZF$
yang tidak berdasar-baik, yakni model $\mathfrak{M}$ sedemikian sehingga
$\dots x_2 \in x_1 \in x_0$.
\end{prob}

Karena model tak standar $\Struct{M}$ dari \olref{thm:non-std} ekuivalen
elementer dengan model standar, sejumlah sifat $\Struct{M}$ dapat
diturunkan. Sisa bagian ini ditujukan pada tugas tersebut, yang akan
memungkinkan kita memperoleh pencirian tepat atas model-model !!{enumerable}
tak standar dari $\Theory{N}$.

\begin{enumerate}
\item Tidak ada anggota $\Domain{M}$ yang $\prec$-kurang dari dirinya
  sendiri: kalimat $\lforall[x][\lnot x < x]$ benar dalam $\Struct{N}$ dan
  karena itu benar pula dalam $\Struct{M}$.
\item Dengan penalaran serupa, kita memperoleh bahwa $\prec$ merupakan
  \emph{pengurutan linear} atas $\Domain{M}$, yakni relasi total, irefleksif,
  dan transitif pada $\Domain{M}$.
\item Unsur $\mathbf{z}$ adalah unsur $\prec$-terkecil dari $\Domain{M}$.
\item Setiap anggota $\Domain{M}$ $\prec$-kurang dari penerus-$*$-nya, dan
  $x^*$ adalah anggota $\prec$-terkecil dari $\Domain{M}$ yang lebih besar
  daripada $x$.
\item $\Struct{M}$ memuat suatu segmen awal (dari $\prec$) yang isomorfik
  dengan $\Nat$: $\mathbf{z}, \mathbf{z}^*, \mathbf{z}^{**}, \dots$, yang
  kita sebut \emph{bagian standar} dari $\Domain{M}$. Setiap anggota lain
  dari $\Domain{M}$ bersifat \emph{tak standar}. Harus terdapat anggota tak
  standar dari $\Domain{M}$; jika tidak, fungsi $h$ dari bukti
  \olref{thm:non-std} merupakan isomorfisme. Kita menggunakan $n, m, \dots$
  sebagai !!{variable}s yang menjangkau bagian standar $\Struct{M}$ ini.
\item Setiap unsur tak standar lebih besar daripada setiap unsur standar;
  sebab, untuk setiap $n \in \Nat$,
  \[
  \Sat{N}{\lforall[z][(\lnot(\eq[z][\Obj{0}] \lor
  \dots \lor \eq[z][\num{n}]) \lif \num{n} < z)]},
  \]
  sehingga jika $z \in \Domain{M}$ berbeda dari semua unsur standar, maka
  unsur itu harus \emph{lebih besar} daripada semuanya.
\item Setiap anggota $\Domain{M}$ selain $\mathbf{z}$ merupakan
  penerus-$*$ dari suatu unsur tunggal $\Domain{M}$, yang dilambangkan
  dengan $^*x$. Jika $x = y^*$, maka $x$ dan $y$ keduanya standar apabila
  salah satunya standar (dan keduanya tak standar apabila salah satunya tak
  standar).
\item Definisikan relasi ekuivalensi $\approx$ pada $\Domain{M}$ dengan
  menyatakan bahwa $x \approx y$ jika dan hanya jika, untuk suatu $n$
  \emph{standar}, berlaku $x \oplus n = y$ atau $y \oplus n =x$. Dengan
  kata lain, $x \approx y$ jika dan hanya jika $x$ dan $y$ terpisah oleh
  jarak berhingga. Jika $n$ dan $m$ standar, maka $n \approx m$. Definisikan
  \emph{blok} dari $x$ sebagai kelas ekuivalensi $[x] = \Setabs{y}{x
  \approx y}$.
\item Andaikan $x \prec y$ dengan $x \not\approx y$. Karena
  $\Sat{N}{\lforall[x][\lforall[y][(x < y \lif (x' < y \lor x' =
      y))]]}$, berlaku $x^* \prec y$ atau $x^* = y$. Yang terakhir mustahil
  karena menyiratkan $x \approx y$, sehingga $x^* \prec y$. Serupa dengan
  itu, jika $x \prec y$ dan $x \not\approx y$, maka $x \prec {^*y}$.
  Karena itu, jika $x \prec y$ dan $x \not\approx y$, maka setiap
  $w \approx x$ $\prec$-kurang dari setiap $v \approx y$. Dengan demikian,
  setiap blok $[x]$ membentuk rantai tak hingga ganda
  \[
  \cdots \prec {^{**}x} \prec {^*}x \prec x \prec x^* \prec x^{**}
  \prec \cdots
  \]
  yang disebut rantai-$Z$ karena memiliki tipe urutan bilangan bulat.
\item Pengurutan $\prec$ dapat diangkat ke blok-blok: jika $x \prec y$,
  maka blok $x$ lebih kecil daripada blok $y$. Suatu blok bersifat
  \emph{tak standar} jika memuat unsur tak standar. Blok standar adalah
  blok terkecil.
\item Tidak ada blok tak standar terkecil: jika $y$ tak standar, terdapat
  $x \prec y$ dengan $x$ juga tak standar dan $x \not\approx y$. Bukti:
  dalam model standar $\Struct{N}$, setiap bilangan habis dibagi dua,
  mungkin dengan sisa satu, yakni
  $\Sat{N}{\lforall[y][\lexists[x][(y = x + x \lor y = x + x +
        \Obj{0}')]]}$. Menurut ekuivalensi elementer, untuk setiap $y \in
  \Domain{M}$ terdapat $x \in \Domain{M}$ sedemikian sehingga berlaku
  $x \oplus x = y$ atau $x \oplus x \oplus \mathbf{z}^*= y$. Jika $x$
  standar, maka $y$ pun standar; jadi $x$ tak standar. Selain itu, $x$ dan
  $y$ termasuk dalam blok yang berbeda, yakni $x \not\approx y$. Untuk
  melihatnya, andaikan keduanya termasuk dalam blok yang sama, yakni
  $x \oplus n = y$ untuk suatu $n$ standar. Jika $y = x \oplus x$, maka
  $x \oplus n = x \oplus x$, sehingga $x = n$ menurut hukum pembatalan
  untuk penjumlahan (yang berlaku dalam $\Struct{N}$ dan karena itu juga
  dalam $\Struct{M}$), dan ternyata $x$ standar. Demikian pula jika
  $y = x \oplus x \oplus \mathbf{z}^*$.
\item Dengan argumen serupa, tidak ada blok terbesar.
\item Pengurutan blok-blok tersebut rapat: jika $[x]$ lebih kecil daripada
  $[y]$ (dengan $x \not\approx y$), maka terdapat blok $[z]$ yang berbeda
  dari keduanya dan terletak di antaranya. Andaikan $x \prec y$. Seperti
  sebelumnya, $x \oplus y$ habis dibagi dua (mungkin dengan sisa), sehingga
  terdapat $u \in \Domain{M}$ sedemikian sehingga berlaku $x \oplus y =
  u \oplus u$ atau $x \oplus y = u \oplus u \oplus \mathbf{z}^*$. Unsur
  $u$ adalah nilai tengah $x$ dan $y$, sehingga terletak di antara keduanya.
  Andaikan $x \oplus y = u \oplus u$ (kasus lain serupa): jika
  $u \approx x$, maka untuk suatu $n$ standar,
  \[
  x \oplus y = x \oplus n \oplus x \oplus n,
  \]
  sehingga $y = x \oplus n \oplus n$ dan akan berlaku $x \approx y$,
  berlawanan dengan asumsi. Kita menyimpulkan bahwa $u \not\approx x$.
  Argumen serupa memberikan $u \not\approx y$.
\end{enumerate}
Dengan demikian, blok-blok tak standar terurut seperti bilangan rasional:
blok-blok tersebut membentuk !!a{enumerable} pengurutan linear tanpa titik
ujung. Akibatnya, untuk setiap dua model !!{enumerable} tak standar
$\Struct{M}_1$ dan $\Struct{M_2}$ dari aritmetika benar, reduk keduanya ke
bahasa yang hanya memuat $<$ dan $=$ adalah isomorfik. Memang, suatu
isomorfisme $h$ dapat didefinisikan sebagai berikut: bagian standar dari
$\Struct{M_1}$ dan $\Struct{M_2}$ isomorfik dengan model standar
$\Struct{N}$, dan karenanya satu sama lain. Blok-blok yang membentuk bagian
tak standar sendiri terurut seperti bilangan rasional dan oleh karena itu,
menurut \olref[dlo]{thm:cantorQ}, isomorfik; suatu isomorfisme antarblok
dapat diperluas menjadi isomorfisme \emph{di dalam} blok dengan memasangkan
unsur-unsur sebarang dalam masing-masing blok, kemudian menetapkan citra
penerus $x$ dalam $\Struct{M_1}$ sebagai penerus citra $x$ dalam
$\Struct{M_2}$. Perhatikan bahwa \emph{tidak} berarti $\mathfrak{M}_1$ dan
$\mathfrak{M}_2$ isomorfik dalam bahasa penuh aritmetika (memang,
isomorfisme selalu relatif terhadap suatu tanda tangan), karena penjumlahan
dan perkalian pada $\Domain{M_1}$ dan $\Domain{M_2}$ dapat didefinisikan
dengan cara-cara yang tidak isomorfik. (Hal ini juga mengikuti dari teorema
terkenal Vaught bahwa banyaknya model terhitung dari suatu teori lengkap
tidak mungkin~2.)

\begin{prob}
Tunjukkan bahwa tidak mungkin ada blok terbesar dalam model aritmetika tak
standar.
\end{prob}

\begin{prob}
Misalkan $\Lang{L}$ adalah !!{language} orde pertama yang hanya memuat $<$
sebagai !!{predicate} (selain $\eq$), dan misalkan $\Struct{N} = (\Nat,
<)$. Semua himpunan bagian berhingga atau kofinit dari $\Struct{N}$ dapat
didefinisikan. Tunjukkan bahwa \emph{hanya} himpunan-himpunan bagian inilah
yang dapat didefinisikan dalam $\Struct{N}$.

(Petunjuk: Pertama, misalkan $\Obj{prc}(x,y)$ adalah rumus $\Lang{L}$ yang
menyingkat ``$x$ adalah pendahulu langsung dari $y$:''
\[
x<y \land \lnot \lexists[z][(x<z \land z < y)].
\]
Sekarang, setiap himpunan bagian $\Struct{N}$ yang dapat didefinisikan
bersesuaian dengan rumus $!A(x)$ dalam $\Lang{L}$. Untuk setiap $!A$ semacam
itu, perhatikan kalimat $!D$:
\[
\lexists[x][\lforall[y][\lforall[z][((x<y \land x<z \land \Obj{prc}(y,z)
  \land !A(y)) \lif !A(z))]]].
\]
Tunjukkan bahwa $\Sat{N}{!D}$ jika dan hanya jika himpunan bagian
$\Struct{N}$ yang didefinisikan oleh $!A$ berhingga atau kofinit.

Sekarang, misalkan $\Struct{M}$ adalah model tak standar yang ekuivalen
elementer dengan $\Struct{N}$. Jika $a \in \Domain{M}$ tak standar,
misalkan $b, c \in \Domain{M}$ lebih besar daripada $a$, dan misalkan $b$
adalah pendahulu langsung dari~$c$. Maka terdapat automorfisme $h$ dari
$\Domain{M}$ sedemikian sehingga $h(b)=c$ (mengapa?). Karena itu, jika $b$
memenuhi $!A$, maka $c$ juga memenuhinya (mengapa?). Akibatnya, $!D$ benar
dalam $\Struct{M}$, dan karenanya juga dalam $\Struct{N}$. Namun, ini
menyiratkan bahwa himpunan bagian $\Struct{N}$ yang didefinisikan oleh $!A$
berhingga atau kofinit.
\end{prob}

\end{document}
